\chapter{The Fisher Information Metric}

The final module of the Schuller curriculum examines the ``geometry of probability'' itself. This field, known as \emph{information geometry}, introduces a Riemannian metric on the space of parameters of a statistical model (the \emph{statistical manifold}), allowing for a geometric interpretation of fluctuations, response functions, and phase transitions.

\section{The Statistical Manifold}

\begin{definition}[Statistical manifold]
A \textbf{statistical manifold} is a smooth manifold $\M$ whose points are probability distributions. A coordinate system $\theta = (\theta^1, \ldots, \theta^m)$ on $\M$ parameterises a family of distributions $\{p(x; \theta)\}_{\theta \in \M}$, where $x$ denotes the random variable (or microstate).
\end{definition}

In thermodynamics, the coordinates on the statistical manifold are typically the intensive variables. For the canonical ensemble:
\begin{equation}
\theta = (\beta, \beta P, -\beta\mu_{\text{chem}}, \ldots),
\end{equation}
and the corresponding probability distribution is the Gibbs measure:
\begin{equation}
p(x; \theta) = \frac{1}{\Partition(\theta)}\exp\!\left(-\sum_i \theta^i\,F_i(x)\right),
\end{equation}
where $F_i(x)$ are the conserved quantities (energy, volume, particle number, etc.). This is an \emph{exponential family} in the language of statistics.

\section{The Fisher Information Metric}

\begin{definition}[Fisher information metric]
The \textbf{Fisher information metric} on the statistical manifold $\M$ is the Riemannian metric $g$ with components:
\begin{equation}
\boxed{g_{ij}(\theta) = \int p(x; \theta)\,\pdv{\ln p(x;\theta)}{\theta^i}\,\pdv{\ln p(x;\theta)}{\theta^j}\,\dd x = \left\langle \pdv{\ell}{\theta^i}\,\pdv{\ell}{\theta^j} \right\rangle,}
\end{equation}
where $\ell(x; \theta) = \ln p(x; \theta)$ is the log-likelihood function.
\end{definition}

The Fisher metric quantifies the ``distance'' between two nearby probability distributions based on their \emph{statistical distinguishability}. If two distributions at $\theta$ and $\theta + \dd\theta$ are easily distinguished by observations, the Fisher distance $\dd s^2 = g_{ij}\,\dd\theta^i\,\dd\theta^j$ is large; if they are nearly identical, the distance is small.

\begin{remark}
The Fisher metric is the unique (up to a constant factor) Riemannian metric on the space of probability distributions that is invariant under sufficient statistics --- this is the content of the \v{C}encov (Chentsov) uniqueness theorem. In this sense, the Fisher metric is the ``natural'' geometry of probability.
\end{remark}

\subsection{Equivalent form via the Hessian}

An equivalent expression for the Fisher metric is:
\begin{equation}
g_{ij} = -\left\langle \frac{\partial^2 \ell}{\partial\theta^i\,\partial\theta^j} \right\rangle.
\end{equation}

This form makes manifest that the Fisher metric measures the \emph{curvature} (in the sense of second derivatives) of the log-likelihood --- distributions that are ``sharply peaked'' in parameter space correspond to large Fisher information.

\section{The Fisher Metric for Exponential Families}

For an exponential family with natural parameters $\theta^i$ and partition function $\Partition(\theta)$:
\begin{equation}
p(x; \theta) = \exp\!\left(\sum_i \theta^i F_i(x) - \Psi(\theta)\right),
\end{equation}
where $\Psi(\theta) = \ln\Partition(\theta)$ is the \textbf{Massieu potential} (log-partition function), the Fisher metric simplifies dramatically:
\begin{equation}
\boxed{g_{ij} = \frac{\partial^2 \Psi}{\partial\theta^i\,\partial\theta^j}.}
\end{equation}

The Fisher metric is the \emph{Hessian of the Massieu potential}. Since $\Psi$ is convex for exponential families, the Hessian is positive-definite, confirming that $g$ is a bona fide Riemannian metric.

\section{Thermodynamic Response Functions as Metric Components}

The components of the Fisher metric are directly related to the thermodynamic response functions. For the canonical ensemble with $\theta = (\beta)$:
\begin{equation}
g_{\beta\beta} = \pdv[2]{\Psi}{\beta} = -\pdv{\langle E \rangle}{\beta} = k_B T^2\,C_V,
\end{equation}
where $C_V$ is the \textbf{heat capacity at constant volume}. More generally, for a system with $\theta = (\beta, \beta P)$:
\begin{equation}
g_{ij} = \begin{pmatrix}
k_B T^2\,C_V & \text{cross-response} \\
\text{cross-response} & k_B T\,\kappa_T\,V
\end{pmatrix},
\end{equation}
where $\kappa_T$ is the \textbf{isothermal compressibility}. The off-diagonal terms are related to cross-susceptibilities.

\begin{proposition}
The Fisher information metric identifies the thermodynamic response functions --- heat capacity, compressibility, susceptibility --- as the \textbf{geometric ``lengths''} of the statistical manifold. A large response function corresponds to a large metric component, meaning that small changes in the corresponding intensive variable lead to large changes in the probability distribution.
\end{proposition}

\section{The Fisher Metric as a Covariance Matrix}

For an exponential family, the Fisher metric components are:
\begin{equation}
g_{ij} = \pdv[2]{\Psi}{\theta^i}{\theta^j} = \langle F_i F_j \rangle - \langle F_i \rangle\langle F_j \rangle = \text{Cov}(F_i, F_j).
\end{equation}

The Fisher metric is the \textbf{covariance matrix of the fluctuations} in the extensive variables. This provides the deep connection between information geometry and thermodynamic fluctuation theory:

\begin{table}[H]
\centering
\renewcommand{\arraystretch}{1.6}
\begin{tabular}{ll}
\toprule
\textbf{Geometric Feature} & \textbf{Physical Interpretation} \\
\midrule
Metric $g_{ij}$ & Covariance of fluctuations (e.g.\ $\langle \Delta E^2\rangle$, $\langle \Delta V^2\rangle$) \\
Geodesic distance & Informational difference between equilibrium states \\
Large $g_{ij}$ & High sensitivity to parameter changes \\
\bottomrule
\end{tabular}
\end{table}

\section{Dual Connections and the $\alpha$-Geometry}

The statistical manifold carries a richer geometric structure than a simple Riemannian manifold. In addition to the Levi-Civita connection $\nabla^{(0)}$ of the Fisher metric, there exists a one-parameter family of affine connections $\nabla^{(\alpha)}$ indexed by $\alpha \in \R$:

\begin{definition}[$\alpha$-connection]
The \textbf{$\alpha$-connection} $\nabla^{(\alpha)}$ is defined by:
\begin{equation}
\Gamma_{ij,k}^{(\alpha)} = \int p(x;\theta)\left(\partial_i\partial_j\ell + \frac{1-\alpha}{2}\,\partial_i\ell\,\partial_j\ell\right)\partial_k\ell\;\dd x,
\end{equation}
where $\partial_i = \partial/\partial\theta^i$ and $\ell = \ln p$.
\end{definition}

The cases $\alpha = \pm 1$ are of particular importance:
\begin{itemize}
\item $\alpha = 1$: the \textbf{exponential connection} (``e-connection''), which is flat for exponential families.
\item $\alpha = -1$: the \textbf{mixture connection} (``m-connection''), which is flat for mixture families.
\item $\alpha = 0$: the Levi-Civita connection of the Fisher metric.
\end{itemize}

The pair $(\nabla^{(1)}, \nabla^{(-1)})$ forms a \textbf{dually flat structure}: the e-connection and m-connection are dual with respect to the Fisher metric. This duality underlies the Legendre transform structure of thermodynamics when viewed from the information-geometric perspective.

\section{Divergence Functions}

The ``distance'' between two points on the statistical manifold can also be measured by divergence functions, which are not symmetric and hence not true metrics:

\begin{definition}[Kullback--Leibler divergence]
The \textbf{KL divergence} from $p(x; \theta)$ to $p(x; \theta')$ is:
\begin{equation}
D_{\text{KL}}(\theta \| \theta') = \int p(x; \theta)\,\ln\frac{p(x;\theta)}{p(x;\theta')}\,\dd x.
\end{equation}
\end{definition}

For nearby distributions ($\theta' = \theta + \dd\theta$), the KL divergence reduces to the Fisher metric:
\begin{equation}
D_{\text{KL}}(\theta \| \theta + \dd\theta) = \frac{1}{2}\,g_{ij}\,\dd\theta^i\,\dd\theta^j + O(\dd\theta^3).
\end{equation}

The Fisher metric is thus the \emph{infinitesimal} form of the KL divergence. In thermodynamics, the KL divergence between two equilibrium states measures the \textbf{irreversible work} (or dissipated free energy) required to transform one state into the other.

\section{Summary}

The Fisher information metric equips the space of equilibrium states with a natural Riemannian geometry:
\begin{itemize}
\item The metric components are the Hessian of the Massieu potential (or free energy).
\item They equal the covariance matrix of fluctuations in the extensive variables.
\item Thermodynamic response functions (heat capacity, compressibility) are geometric invariants.
\item The statistical manifold carries a dually flat structure reflecting the Legendre transform.
\item The KL divergence provides a measure of irreversibility.
\end{itemize}

This sets the stage for the study of curvature and its divergence at phase transitions in the next chapter.
